% Copyright (c) 2026 Geovana Neves. Licensed under CC BY 4.0 (https://creativecommons.org/licenses/by/4.0/). See LICENSE-AND-AUTHORSHIP.md.
%
\section{What every product states}

\begin{frame}{The condition block: angles and velocities}
\small Every table the post stage composes states its condition in one block,
in the order shown across these two frames.
\par\vspace{3mm}
{\ftstablesize
\begin{tabular}{@{}p{0.22\linewidth}p{0.10\linewidth}p{0.59\linewidth}@{}}
\toprule
\textbf{Column} & \textbf{Unit} & \textbf{Meaning} \\
\midrule
\code{ALPHA}, \code{BETA} & deg & Angles the solver reports it ran at, as the row wrote them \\
\code{MACH} & - & Mach number of that point \\
\code{RE} & millions & Reynolds number the solver reports \\
\code{VINF} & m/s & Free-stream velocity the solver reports \\
\code{VREF} & m/s & Solver reference velocity, used to normalise a coefficient \\
\bottomrule
\end{tabular}}
\end{frame}

\begin{frame}{The condition block: atmosphere and references}
\relax
{\ftstablesize
\begin{tabular}{@{}p{0.22\linewidth}p{0.10\linewidth}p{0.59\linewidth}@{}}
\toprule
\textbf{Column} & \textbf{Unit} & \textbf{Meaning, continuing the block} \\
\midrule
\code{ALT} & ft & Altitude the row states; \code{NA} where it states none \\
\code{RHO} & kg/m$^3$ & Air density the run resolved for that point \\
\code{TEMP} & K & Air temperature the run resolved for that point \\
\code{MU} & Pa s & Dynamic viscosity, written in scientific notation \\
\code{J} & - & Advance ratio the row requested; \code{NA} on a row that turns no rotor \\
\code{J\_CLOCK} & - & Advance ratio the clock rotor ran at, $V/(nD)$; \code{NA} where the record or the reference does not say \\
\code{RPM\_CLOCK} & rev/min & Speed the clock rotor turned at, with its hand; \code{NA} with no clock \\
\code{SREF}, \code{CREF}, \code{BREF} & m$^2$, m, m & Reference area, chord and span \\
\bottomrule
\end{tabular}}
\end{frame}

\begin{frame}{A coefficient needs its normalisation}
\small A coefficient is a force divided by
\[
  \frac{1}{2}\rho V^2 S.
\]
\begin{itemize}
  \item Area alone does not let a reader recover a force or compare campaigns.
        Density and the velocity used in the division must be stated too.
  \item \code{ALT} cannot stand in for density: a row may pin density directly.
\end{itemize}
\end{frame}

\begin{frame}{Which velocity normalises which product}
\begin{itemize}\small
  \item The solver uses \code{VREF}. The steady polar's twenty-four
        coefficients are the solver's own and use \code{VREF}.
  \item The plots table, every reduction of it and the unsteady polar are
        rescaled to the free stream by
        $\left(\mathrm{VREF}/\mathrm{VINF}\right)^2$ and use \code{VINF}.
  \item The rotor table recovers Newtons with \code{VREF}, undoing the
        solver's division. Its coefficients use $\rho n^2 D^4$ and carry
        no velocity.
\end{itemize}
\begin{takeaway}
Both velocities are stated. When they differ, post warns and names the point.
\end{takeaway}
\end{frame}

\begin{frame}{Reference area and chord must match the export}
\begin{itemize}\small
  \item The solver divides by the area and length in the project it opened.
        The loads export prints both.
  \item The package sets neither. It states the reference artifact's
        \code{SREF} and \code{CREF}, and checks them against the export.
  \item A difference beyond the export's printed precision is a
        \textbf{warning} in \code{post.log} naming both numbers, and every
        computable product is still written;
        \code{--check-frozen} refuses the simulation's products instead.
\end{itemize}
\begin{takeaway}
A reference beside a coefficient must be the reference that divided it.
\end{takeaway}
\end{frame}

\begin{frame}{Where the condition block appears}
\begin{itemize}
  \item The steady polar already carries \code{ALPHA}, \code{BETA},
        \code{MACH} and \code{RE} among its twenty-four columns. It adds
        the rest beside them.
  \item \code{probes/<point>\_plots.csv} states \textbf{no condition}.
        It keeps the export's own header.
  \item Reductions read every column of that plots table back as a plotted
        quantity.
\end{itemize}
\end{frame}
